\documentclass{article}
\usepackage{amsmath}
\usepackage{amsthm}
\usepackage{amssymb}
\usepackage{setspace}
\theoremstyle{definition}
\newtheorem{definition}{Def}
\newtheorem{theorem}[definition]{Thm}
\newtheorem{corollary}[definition]{Cor}
\newtheorem{proposition}[definition]{Prop}
\newtheorem{remarks}[definition]{Rmk}
\newtheorem{notation}[definition]{Notation}
\onehalfspacing

\title{Chapter2}
\author{Ricky Dai}
\date{September 2026}

\begin{document}

\maketitle

\section{Finite, Countable and Uncountable Sets}
\begin{definition}
    Consider two sets $A$ and $B$, and suppose that with each element of $A$ there is associated, in some manner, an element of $B$, which we denote by $f(x)$. Then $f$ is said to be a \textit{function} from $A$ to $B$ (or a \textit{mapping} of $A$ into $B$). The set $A$ is called the \textit{domain} of $f$ (we also say $f$ is defined on $A$) and the elements $f(x)$ are called the \textit{values} of $f$. The set of all values of $f$ is called the \textit{range} of $f$.
\end{definition}
\begin{definition}
    Let $A$ and $B$ be two sets and let $f$ be a mapping of $A$ into $B$. If $E\subset A$, $f(E)$ is defined to be the set of all elements $f(x)$, for $x\in E$. We call $f(E)$ the \textit{image} of $E$ under $f$. In this notation, $f(A)$ is the range of $f$. It is clear that $f(A)\subset B$. If $f(A)=B$, we say that $f$ maps $A$ onto $B$.
\end{definition}
\begin{definition}
    If there exists a $1-1$ mapping of $A$ \textit{onto} $B$, we say that $A$ and $B$ can be put in $1-1$ \textit{correspondence}, or that $A$ and $B$ have the same \textit{cardinal number}, or, briefly, that $A$ and $B$ are \textit{equivalent}, and we write $A\sim B$. This relation clearly has the following properties:\par
    Reflexive: $A\sim A$.\par
    Symmetric: If $A\sim B$, then $B\sim A$.\par
    Transitive: If $A\sim B$ and $B\sim C$, then $A\sim C$.\par
    Any relation with these three properties is called an \textit{equivalence relation}.
\end{definition}
\begin{definition}
    For any positive $n$, let $\mathbb{N}_n$ be the set whose elements are the integers $1, 2, ..., n$; let $\mathbb{N}$ be the set consisting of all positive integers. For any set $A$, we say:\par
    (a) $A$ is finite if $A\sim \mathbb{N}_n$ for some $n$ (the empty set is also considered to be finite).\par
    (b) $A$ is \textit{infinite} if $A$ is not finite.\par
    (c) $A$ is \textit{countable} if $A\sim \mathbb{N}$.\par
    (d) $A$ is \textit{uncountable} if $A$ is neither finite nor countable.\par
    (e) $A$ is \textit{at most countable} if $A$ is finite or countable.\par
    Countable sets are sometimes called \textit{enumerable}, or \textit{denumerable}.
\end{definition}
\begin{remarks}
    In the definition above, (b) can be replaced by the statement: $A$ is infinite if $A$ is equivalent to one of its proper subsets.
\end{remarks}
\begin{definition}
    By a \textit{sequence}, we mean a function $f$ defined on the set $\mathbb{N}$ of all positive integers. If $f(n)=x_n$, for $n\in \mathbb{N}$, it is customary to denote the sequence $f$ by the symbol $(x_n)_{n=1}^\infty$, or more briefly by $(x_n)$, or sometimes by $x_1, x_2, x_3,...$. The values of $f$, that is, the elements $x_n$, are called the \textit{terms} of the sequence. If $A$ is a set and if $x_n\in A$ for all $n \in \mathbb{N}$, then $(x_n)$ is said to be a \textit{sequence in} $A$, or a \textit{sequence of elements of} $A$.
\end{definition}
\begin{theorem}
    Every infinite subset of a countable set $A$ is countable.
\end{theorem}
\begin{definition}
    Let $A$ and $\Omega$ be sets, and suppose that with each element $\alpha$ of $A$ there is associated a subset of $\Omega$ which we denote by $E_\alpha$.\par
    The set whose elements are the sets $E_\alpha$ will be denoted by $\{E_\alpha\}$.
\end{definition}
\begin{remarks}
    communicative law: $A\bigcup B = B\bigcup A$; $A\bigcap B = B\bigcap A$.\par
    associative law: $(A\bigcup B)\bigcup C = A \bigcup(B\bigcup C)$; $(A\bigcap B)\bigcap C = A\bigcap (B\bigcap C)$.\par
    distribution law: $A\bigcap(B\bigcup C) = (A\bigcap B)\bigcup(A\bigcap C)$.
\end{remarks}
\begin{theorem}
    Let $\{E_n\}$, $n = 1,2,3,...,$ be a countable collection of countable sets, and put
    \[
    S=\bigcup_{n=1}^\infty E_n.
    \]
    Then S is countable.
\end{theorem}
\begin{corollary}
    Suppose $A$ is at most countable, and, for every $\alpha\in A$, $B_\alpha$ is at most countable. Put
    \[
    T=\bigcup_{\alpha\in A} B_\alpha.
    \]
    Then $T$ is at most countable, for $T$ is equivalent to a subset of S as stated above.
\end{corollary}
\begin{theorem}
    Let $A$ be a countable set, and let $B_n$ be the set of all $n-tuples(a_1,..., a_n)$, where $a_k\in A\;(k=1,...,n)$, and the elements $a_1,...,a_n$ need not be distinct. Then $B_n$ is countable.
\end{theorem}
\begin{corollary}
    The set of all rational numbers is countable.
\end{corollary}
\begin{theorem}
    Let $A$ be the set of all sequences whose elements are the digits $0$ and $1$. This set $A$ is uncountable.
\end{theorem}

\section{Metric Spaces}
\begin{definition}
    A set $X$, whose elements we shall call \textit{points}, is said to be a \textit{metric space} if with any two points $p$ and $q$ of $X$ there is associated a real number $d(p, q)$, called the \textit{distance} from $p$ to $q$, such that\par
    (a) $d(p, q)>0$ if $p\neq q$; $d(p, p)=0$;\par
    (b) $d(p, q)=d(q, p)$;\par
    (c) $d(p, q) \leq d(p, r) + d(r, q)$ for any $r\in X$.\\
    Any function with these three properties is called a \textit{distance function}, or a \textit{metric}.
\end{definition}
\begin{definition}
    By the \textit{segment} $(a, b)$ we mean the set of all real numbers $x$ such that $a<x<b$.\par
    By the \textit{interval} $[a,b]$ we mean the set of all real numbers $x$ such that $a\leq x\leq b$.\par
    If $a_i < b_i$ for $i=1, ..., k$, the set of all points $\textbf{x}=(x_1,...x_k)$ in $\mathbb{R}^k$ whose coordinates satisfy the inequalities $a_i\leq x\leq b_i\;(1\leq i\leq k)$ is called a \textit{k-cell}. Thus a 1-cell is an interval, a 2-cell is a rectangle, etc.\par
    If $\mathbf{x}\in \mathbb{R}^k$ and $r>0$, the \textit{open (or closed) ball} $B$ with center at $\mathbf{x}$ and radius $r$ is defined to be the set of all $\mathbf{y}\in \mathbb{R}^k$ such that $|\mathbf{y}-\mathbf{x}|<r$ (or $|\mathbf{y}-\mathbf{x}|\leq r$).\par
    We call a set $E\subset \mathbb{R}^k$ \textit{convex} if
    \[
    \lambda\mathbf{x}+(1-\lambda)\mathbf{y}\in E
    \]
    whenever $\mathbf{x}\in E,\; \mathbf{y}\in E$, and $0<\lambda<1$.
\end{definition}
\begin{definition}
    Let $X$ be a metric space. All points and sets mentioned below are understood to be elements and subsets of $X$.\par
    (a) \textit{A neighborhood} of $p$ is a set $N_r(p)$ consisting of all $q$ such that $d(p, q)<r$ for some $r>0$. The number $r$ is called the \textit{radius} of $N_r(p)$.\par
    (b) A point $p$ is a \textit{limit point} of the set $E$ if \textit{every} neighborhood of $p$ contains a point $q\neq p$ such that $q\in E$.\par
    (c) If $p\in E$ and $p$ is not a limit point of $E$, then $p$ is called an \textit{isolated point} of $E$.\par
    (d) $E$ is \textit{closed} if every limit point of $E$ is a point of $E$.\par
    (e) A point $p$ is an \textit{interior} point of $E$ if there is a neighborhood $N$ of $p$ such that $N\subset E$.\par
    (f) $E$ is \textit{open} if every point of $E$ is an interior point of $E$.\par
    (g) The \textit{complement} of $E$ (denoted by $E^c$) is the set of all points $p\in X$ such that $p\notin E$.\par
    (h) $E$ is \textit{perfect} if $E$ is closed and if every point of $E$ is a limit point of $E$.\par
    (i) $E$ is \textit{bounded} if there is a real number $M$ and a point $q\in X$ such that $d(p, q)<M$ for all $p\in E$.\par
    (j) $E$ is \textit{dense in} $X$ if every point of $X$ is a limit point of $E$, or a point of $E$ (or both). 
\end{definition}
\begin{theorem}
    Every neighborhood is an open set.
\end{theorem}
\begin{theorem}
    If $p$ is a limit point if a set $E$, then every neighborhood of $p$ contains infinitely many points of $E$.
\end{theorem}
\begin{corollary}
    A finite point set has no limit points.
\end{corollary}
\begin{theorem}
    Let $\{E_\alpha\}$ be a (finite or infinite) collection of sets $E_\alpha$. Then
    \[
    (\bigcup_\alpha E_\alpha)^c=\bigcap_\alpha \;(E_\alpha^c).
    \]
\end{theorem}
\begin{theorem}
    A set $E$ is open if and only if its complement is closed.
\end{theorem}
\begin{corollary}
    A set $F$ is closed if and only if its complement is open.
\end{corollary}
\begin{theorem}
    (a) For any collection $\{G_\alpha\}$ of open sets, $\bigcup_\alpha G_\alpha$ is open.\par
    (b) For any collection $\{F_\alpha\}$ of closed sets, $\bigcap_\alpha F_\alpha$ is closed.\par
    (c) For any finite collection $G_1,...,G_n$ of open sets, $\bigcup_{i=1}^n G_i$ is open.\par
    (d) For any finite collection $F_1,...,F_n$ of closed sets, $\bigcap_{i=1}^n F_i$ is closed.
\end{theorem}
\begin{definition}
    If $X$ is a metric space, if $E\subset X$, and if $E'$ denotes the set of all limit points of $E$ in $X$, then the \textit{closure} of $E$ is the set $\overline{E}=E\bigcup E'$.
\end{definition}
\begin{theorem}
    If $X$ is a metric space and $E\subset X$, then\par
    (a) $\overline{E}$ is closed,\par
    (b) $E=\overline{E}$ if and only if $E$ is closed,\par
    (c) $\overline{E}\subset F$ for every closed set $F\subset X$ such that $E\subset F$.
\end{theorem}
\begin{theorem}
    Let $E$ be a nonempty set of real numbers which is bounded above. Let $y=supE$. Then $y\in \overline{E}$. Hence $y\in E$ if $E$ is closed.
\end{theorem}
\begin{theorem}
    Suppose $Y\subset X$. A subset $E$ of $Y$ is open relative to $Y$ if and only if $E=Y\bigcap G$ for some open subset $G$ of $X$.
\end{theorem}
\section{Compact Sets}
\begin{definition}
    By an \textit{open cover} of a set $E$ in a metric space $X$ we mean a collection $\{G_\alpha\}$ of open subset of $X$ such that $E\subset \bigcup_\alpha G_\alpha$.
\end{definition}
\begin{definition}
    A subset $K$ of a metric space $X$ is said to be \textit{compact} if every open cover of $K$ contains a \textit{finite} subcover.
\end{definition}
\begin{theorem}
    Suppose $K\subset Y\subset X$. Then $K$ is compact relative to $X$ if and only if $K$ is compact relative to Y.
\end{theorem}
\begin{theorem}
    Compact subsets of metric spaces are closed.
\end{theorem}
\begin{theorem}
    Closed subsets of compact sets are compact.
\end{theorem}
\begin{corollary}
    If $F$ is closed and $K$ is compact, then $F\bigcap K$ is compact.
\end{corollary}
\begin{theorem}
    If $\{K_\alpha\}$ is a collection of compact subsets of a metric space $X$ such that the intersection of every finite subcollection of $\{K_\alpha\}$ is nonempty, then $\bigcap K_\alpha$ is nonempty.
\end{theorem}
\begin{corollary}
    If $\{K_n\}$ is a countable collection of nonempty compact sets such that $K_n\supset K_{n+1}\;(n=1,2,3,...),$ then $\bigcap_{n=1}^\infty K_n$ is not empty
\end{corollary}
\begin{theorem}
    If $E$ is an infinite subset of a compact set $K$, then $E$ has a limit point in $K$.
\end{theorem}
\begin{theorem}
    If $\{I_n\}$ is a sequence of intervals in $\mathbb{R}^1$, such that $I_n \supset I_{n+1}\;(n=1,2,3,... ),$ then $\bigcap_{n=1}^\infty I_n$ is not empty.
\end{theorem}
\begin{theorem}
    Let $k$ be a positive integer. If $\{I_n\}$ is a sequence of k-cells such that $I_n\supset I_{n+1}\;(n=1,2,3,...)$, then $\bigcap_{n=1}^\infty I_n$ is not empty.
\end{theorem}
\begin{theorem}
    Every k-cell is compact.
\end{theorem}
\begin{theorem}
    If a set $E$ in $\mathbb{R}^k$ has one of the following three properties, then it has the other two:\par
    (a) $E$ is closed and bounded.\par
    (b) $E$ is compact.\par
    (c) Every infinite subset of $E$ has a limit point in $E$.
\end{theorem}
\begin{theorem}
    (Weierstrass). Every bounded infinite subset of $\mathbb{R}^k$ has a limit point in $\mathbb{R}^k$.
\end{theorem}
\section{Perfect Sets}
\begin{theorem}
    Let $P$ be a nonempty perfect set in $\mathbb{R}^k$. Then $P$ is uncountable.
\end{theorem}
\begin{corollary}
    Every interval $[a, b]\;(a<b)$ is uncountable. In particular, the set of all real numbers is uncountable.
\end{corollary}
\begin{definition}
    \textbf{The Cantor Set.} The set which we are now going to construct shows that there exist perfect sets in $\textit{R}^1$ which contain no segment.\par
    Let $E_0$ be the interval $[0, 1]$. Remove the segment $(1/3, 2/3)$, and let $E_1$ be the union of the intervals
    \[
    [0,\frac{1}{3}],[\frac{2}{3},1].
    \]
    Remove the middle thirds of these intervals, and let $E_2$ be the union of the intervals
    \[
    [0,\frac{1}{9}],[\frac{2}{9},\frac{3}{9}],[\frac{6}{9},\frac{7}{9}],[\frac{8}{9},1]
    \]
    Continuing in the way, we obtain a sequence of compact sets $E_n$ such that\par
    (a) $E_1\supset E_2\supset E_3\supset ...$;\par
    (b) $E_n$ is the union of $2^n$ intervals, each of length $3^{-n}$.\par
    \quad The set
    \[
    P=\bigcap_{n=1}^\infty E_n
    \]
    is called the \textit{Cantor set}. $P$ is clearly compact and nonempty.
\end{definition}
\section{Connected Sets}
\begin{definition}
    Two subsets $A$ and $B$ of a metric space $X$ are said to be \textit{separated} if both $A\bigcap\overline{B}$ and $\overline{A}\bigcap B$ are empty, i.e., if no point of $A$ lies in the closure of $B$ and no point of $B$ lies in the closure of $A$.\par
    A set $E\subset X$ os said to be \textit{connected} if $E$ is \textit{not} a union of two nonempty separated sets.
\end{definition}
\begin{remarks}
    Separated sets are of course disjoint, but disjoint sets need not be separated. For example, the interval $[0,1]$ and the segment $(1, 2)$ are \textit{not} separated since 1 is a limit point of $(1, 2)$. However, the segments $(0,1)$ and $(1,2)$ \textit{are} separated.
\end{remarks}
\begin{theorem}
    A subset $E$ of the real line $\mathbb{R}^1$ is connected if and only if it has the following property: If $x\in E$, $y\in E$, and $x<z<y$, then $z\in E$.
\end{theorem}
\section{Exercise Notes}
\begin{theorem}
    All algebraic numbers is countable.
\end{theorem}
\begin{corollary}
    All transcendental numbers is uncountable.
\end{corollary}
\begin{definition}
    Let $E^\circ$ denote the set of all interior points of a set $E$> Then $E^\circ$ is called the interior of $E$.
\end{definition}
\begin{theorem}
    Every connected metric space with at least two points is uncountable.
\end{theorem}
\begin{definition}
    A metric space is called \textit{separable} if it contains a countable dense subset.
\end{definition}
\begin{definition}
    A collection $\{V_\alpha\}$ of open subsets of $X$ is said to be a \textit{base} for $X$ if the following is true: For every $x\in X$ and every open set $G\subset X$ such that $x\in G$, we have $x\in V_\alpha \subset G$ for some $\alpha$. In other words, every open set in $X$ is the union of a subcollection of $\{V_\alpha\}$.
\end{definition}
\begin{theorem}
    Every separable metric space has a \textit{countable} base.
\end{theorem}
\begin{theorem}
    If a metric space has a countable base, then it is separable.
\end{theorem}
\begin{corollary}
    For a metric space, having a countable base $\iff$ being separable.
\end{corollary}
\begin{theorem}
    Every compact metric space $K$ has a countable base, and that $K$ is therefore separable.
\end{theorem}
\begin{theorem}
    Let $X$ be a metric space in which every infinite subset has a limit point. Then $X$ is compact.
\end{theorem}
\begin{definition}
    A point $p$ in a metric space $X$ is a \textit{condensation point} of a set $E\subset X$ if every neighborhood of $p$ contains uncountably many points of $E$.
\end{definition}
\begin{theorem}
    Suppose $E\subset \mathbb{R}^k$, $E$ is uncountable, and let $P$ be the set of all condensation points of $E$. Then $P$ is perfect and that at most countably many points of $E$ are not in $P$.
\end{theorem}
\begin{theorem}
    If $\mathbb{R}^k=\bigcup_{n=1}^\infty F_n$, where each $F_n$ is a closed subset of $\mathbb{R}^k$, then at least one $F_n$ has a nonempty interior.\par
    \textit{Equivalent statement:} If $G_n(=\mathbb{R}^k\setminus F_n)$ is a dense open subset of $\mathbb{R}^k$ for $n\;=\;1,2,3,...,$ then $\bigcap_{n=1}^\infty G_n$ is not empty.
\end{theorem}
\begin{remarks}
    \textbf{Relations between Some Concepts:}\par
    \begin{tabular}{@{}l@{\quad}l@{\quad}l@{}}
    Compact   & $\iff$             & Closed + Bounded $\iff$ Every infinite subset has a limit point in the set; \\
    Compact   & $\Longrightarrow$ & Separable; \\
    Compact   & $\Longrightarrow$ & Every open cover contains a \textit{finite} subcover;\\
    Separable & $\iff$             & Having a countable base; \\
    Perfect   & $\iff$             & Closed + No isolated points (or say every point is a limit point); \\
    Cantor Set& $\Longrightarrow$ & Compact and Perfect.
    \end{tabular}
\end{remarks}
\end{document}
